%% kapa: exam=1 round=27 year=2016 period=1 form=A subject=civil_law subject_ko=민법 seq=39 number=39
%% source: Q-Net 감정평가사 기출문제 (cst00309) artlSeq=5169402 file=extracted/1008/5169402/제27회 감정평가사 1차 문제/제27회 감정평가사 1교시문제(A형).pdf
%% conversion: kordoc 4.24.51 -> tools/gen1.py (auto); review=proofread-sonnet; flags=none
\begin{question}{39}
법정지상권에 관한 설명으로 옳은 것은? (다툼이 있으면 판례에 따름)
\begin{choices}
  \item 법정지상권의 성립 후 구건물이 철거되고 신건물이 축조된 경우, 그 법정지상권의 존속기간ㆍ범위 등은 신건물을 기준으로 한다.
  \item 관습상의 법정지상권에서 건물은 등기가 되어 있지 않아도 무방하나, 무허가건물이어서는 안 된다.
  \item 토지와 함께 공동근저당권이 설정된 건물이 그대로 존속함에도 등기가 멸실되고 등기부가 폐쇄되면, 그 후 경매로 토지와 건물의 소유자가 달라지더라도 법정지상권이 성립할 수 없다.
  \item 구분소유적 공유관계에 있는 자가 자신의 특정 소유가 아닌 부분에 신축한 건물을 제3자에게 양도한 경우에 관습상 법정지상권이 성립한다.
  \item 가압류 후 본압류 및 강제경매가 이루어지는 경우 관습상 법정지상권의 요건으로 ‘토지와 그 지상 건물이 동일인 소유’인지 여부는 가압류의 효력 발생 시를 기준으로 한다.
\end{choices}
\end{question}
\answer{5}
